\DTMsavedate{mydate}{2022-06-30}
\maketitle{Engineering Mathematics 1A}{Applied Mathematics}{\DTMusedate{mydate}}

% ----------------------------------------------------- %
% COURSE INFO
\subsection*{Course Information}\label{course:sma3116}
\vspace{0ex}%
\noindent\parbox{0.5\textwidth}{%
    \noindent\begin{tabular}{@{}ll}
                 \textsf{Course:}          & 	Engineering Mathematics 1A \Romannum{4}   \\
                 \textsf{Course Code:}         & 	SMA 1116    \\
                 \textsf{Credit Hours:}    & 	10
    \end{tabular}}
\vspace{2ex}\hrule\vspace{2ex}

% ----------------------------------------------------- %
% INSTRUCTOR INFO
\subsection*{Lecturer Information}
\noindent\parbox{0.5\textwidth}{%
    \noindent\begin{tabular}{@{}ll}

    \end{tabular}}
\vspace{2ex}\hrule\vspace{2ex}


% ----------------------------------------------------- %
% COURSE SYNOPSISDESCRIPTION
\subsection*{Course Synopsis}
The module explores calculus in one variable i.e.\ Limits and continuity of functions, Differentiation,
Leibniz’s Rule, L’Hopitals Rule, Elementary functions including hyperbolic functions and their
inverses; Integration – techniques including reduction formulae; Applications – arc-length, area,
volumes, moments of inertia, centroids; Plane polar coordinates; Complex Numbers: Basic algebra; De
Moivre’s theorem; Complex exponentials; Linear Algebra: Vector algebra in 2 and 3 dimensions; Scalar
and vector products; Equations of lines and planes.
\vspace{2ex} \hrule\vspace{2ex}

% ----------------------------------------------------- %
% LEARNING OBJECTIVES
\subsection*{Learning Objectives}
On completion of the course students should be able to:
\begin{outline}
    \1

    \1

    \1

    \1

    \1

    \1

    \1

    \1

    \1

    \1

    \1

\end{outline}
\vspace{2ex}\hrule\vspace{2ex}


% ----------------------------------------------------- %
% COURSE TOPICS
\subsection*{Topics}
\begin{outline}
    \1
    \begin{itemize}
        \item
        \item
        \item
        \item
        \item
    \end{itemize}
    \1
    \1
\end{outline}
\vspace{2ex} \hrule\vspace{2ex}

% ----------------------------------------------------- %
% Students Assenssment
\subsection*{Assessment of Students}
\begin{outline}
    \1 Students are assessed through:
    \begin{enumerate}
        \item course work consisting of two in-class tests, test on one Ordinary Differential Equations and Partial Differential Equations and test two on Numerical Methods.
        \item end of semester final examination with Section A of 40 marks and Section B of 60 marks.
    \end{enumerate}

\end{outline}
\vspace{2ex}\hrule\vspace{2ex}

% ----------------------------------------------------- %
% Grade Evaluation
\subsection*{Course Evaluation and Grading Policies}
Grades are based on:
Continuous assessment (\qty{25}{\percent}),
Final Exam (\qty{75}{\percent})
\vspace{2ex}\hrule\vspace{2ex}

% ----------------------------------------------------- %
% EQUIPMENT
% https://sites.udel.edu/computing-purchases/personal-specs/
% \subsection*{Essential Equipment}

% \vspace{2ex}\hrule\vspace{2ex}

% Policies and Other information
\subsection*{Policies and Other Information}
% Conduct
\subsection*{Key Text}
\begin{itemize}
    \item
    \item
    \item
\end{itemize}
\vspace{2ex}\hrule
\vspace{2ex}\hrule\vspace{2ex}